\capitulofrontal{Fontes e créditos}{creditos}

\secaofrontal{Questões oficiais}
As questões identificadas com o selo \chip{ebookDark}{\faUniversity\ ENEM} foram aplicadas no
Exame Nacional do Ensino Médio e são reproduzidas a partir dos cadernos de prova oficiais
(caderno azul do 2.º dia, edições \anosbanco), disponibilizados publicamente pelo INEP em
\url{https://www.gov.br/inep/pt-br/areas-de-atuacao/avaliacao-e-exames-educacionais/enem/provas-e-gabaritos}.
O número indicado em cada questão é o do caderno azul. Os gabaritos foram conferidos com os
gabaritos oficiais publicados pelo INEP.

\secaofrontal{Parâmetros da TRI}
Os parâmetros $a$ (discriminação), $b$ (dificuldade) e $c$ (acerto casual) e a habilidade avaliada
por cada questão foram obtidos dos \emph{Microdados do ENEM} (arquivo de itens da prova) publicados
pelo INEP em \url{https://www.gov.br/inep/pt-br/acesso-a-informacao/dados-abertos/microdados/enem}.
A posição aproximada na escala ($500 + 100b$) e a classificação em quatro níveis são critérios
desta coleção. Itens anulados ou sem parâmetros divulgados não foram utilizados.

\secaofrontal{Matriz de Referência}
INEP. \emph{Matriz de Referência ENEM} — Matemática e suas Tecnologias e Anexo dos objetos de
conhecimento associados, disponível em \url{https://download.inep.gov.br/download/enem/matriz_referencia.pdf}.

\secaofrontal{Conteúdo original}
Os textos teóricos, os exemplos, as dicas, os macetes, as questões inéditas e todas as resoluções
comentadas foram elaborados para esta coleção, com apoio de inteligência artificial, e passaram
por revisão matemática e pedagógica. Encontrou algum erro? Anote a página e a questão: toda
correção ajuda os próximos estudantes.

\secaofrontal{Uso}
Material de estudo independente, sem vínculo com o INEP ou com o Ministério da Educação.
Os direitos sobre os textos-base citados nas questões oficiais pertencem aos respectivos autores,
indicados nas próprias questões.
